\section{Complete proofs for the selected LPO and historical KBO}

\begin{annexitem}{annex:p12.ro.lpo-structural}{thm:p12-ro-lpo-structural}{Every structural case of lexicographic path comparison}
\begin{proof}
All comparisons are defined by recursion on the sum of the two term
sizes, as each subsidiary comparison uses a proper subterm. In
particular every term is greater than each immediate subterm by the
first clause, taking equality at that argument.

First prove transitivity by induction on
\(\lvert s\rvert+\lvert t\rvert+\lvert u\rvert\), assuming
\(s>_{\rm lpo}t>_{\rm lpo}u\). If the first comparison uses a
child \(s_i\geq_{\rm lpo}t\), then either \(s_i=t\), so
\(s_i>_{\rm lpo}u\), or the induction hypothesis gives that
comparison from \(s_i>_{\rm lpo}t>_{\rm lpo}u\). The first
clause therefore gives \(s>_{\rm lpo}u\).

Otherwise write \(s=f(\bar s)\), \(t=g(\bar t)\), where
\(s>_{\rm lpo}t_j\) for every \(j\), and the root/lexicographic
condition holds. If \(t>_{\rm lpo}u\) uses a child
\(t_j\geq_{\rm lpo}u\), equality or induction on the smaller
triple \((s,t_j,u)\) gives \(s>_{\rm lpo}u\).
In the remaining case \(u=h(\bar u)\),
\(t>_{\rm lpo}u_k\) for every \(k\), and the second
root/lexicographic condition holds. Induction on \((s,t,u_k)\)
gives \(s>_{\rm lpo}u_k\) for every \(k\).
If either root comparison is strict, transitivity of the precedence,
together with equality in the other comparison when applicable, gives
\(f>_\Sigma h\). If neither is strict, all three roots coincide.
Let the first unequal coordinates in the two lexicographic comparisons
have status indices \(i,j\). When \(i<j\), the first strict
comparison persists because the second pair is still equal at that
coordinate. When \(j<i\), the second persists because the first
pair is equal there. When \(i=j\), induction on the three proper
argument terms gives transitivity at that coordinate. Earlier
coordinates are equal. Thus \(\bar s\) is lexicographically
greater than \(\bar u\), and the second clause gives
\(s>_{\rm lpo}u\). This covers all transitivity cases.

Prove irreflexivity by induction on term size. A variable has no
outgoing comparison. A putative comparison \(s>_{\rm lpo}s\)
cannot use a strict root precedence or a first unequal coordinate of
an identical tuple. The remaining clause would require a proper child
\(s_i\geq_{\rm lpo}s\). Equality is impossible by size. If
\(s_i>_{\rm lpo}s\), combine it with the already established
\(s>_{\rm lpo}s_i\) and transitivity to obtain
\(s_i>_{\rm lpo}s_i\), contradicting induction. Hence the
relation is irreflexive. Repeated immediate-subterm comparisons and
transitivity give the full proper-subterm property.

For context closure it is enough to replace one argument. Suppose
\(s>_{\rm lpo}t\), and let \(S=g(u_1,\ldots,s,\ldots,u_n)\)
and \(T=g(u_1,\ldots,t,\ldots,u_n)\). Every unchanged target
argument is a proper subterm of \(S\). The changed argument \(t\)
is also smaller than \(S\), by \(S>_{\rm lpo}s>_{\rm lpo}t\).
The first unequal coordinate at the fixed status is the changed one,
where the stipulated strict comparison holds. The second defining
clause therefore gives \(S>_{\rm lpo}T\). Induction on the depth
of a one-hole context proves arbitrary context closure.

For substitution stability, induct on the comparison's pair of term
sizes. In the subterm clause, equality of the witnessing child and
the target remains equality after substitution; a strict child
comparison remains strict by induction. In the root clause, neither
root changes. All required comparisons with target children persist
by induction. If the roots were strictly ordered, that condition is
unchanged. Otherwise the preceding equal argument pairs remain equal,
and induction preserves the strict comparison at the first unequal
coordinate. That coordinate cannot become equal, by irreflexivity of
the preserved strict order. The root clause again applies.

Finally induct on the pair of term sizes to prove variable containment.
In the subterm clause, every variable of the target belongs to the
witnessing child by equality or induction, hence to the source. In
the root clause, every variable of the target occurs in a target child;
the domination comparisons and induction put it in the source. A
nullary target has no variables. This finishes all stated properties.
\end{proof}
\end{annexitem}

\begin{annexitem}{annex:p12.ro.finite-tree}{thm:p12-ro-finite-tree}{Finite-ranked tree embedding and the simplification-order argument}
\begin{proof}
Call a sequence bad for a relation \(\preceq\) when it has no
indices \(i<j\) with \(a_i\preceq a_j\). A reflexive transitive
relation is a well-quasi-order when it admits no infinite bad sequence.
The embedding relation is reflexive by induction on terms. It is
transitive by induction on the combined size of three terms: if the
second embedding dives into a child of its target, compose into that
child by induction and dive again. Otherwise it couples equal roots.
If the first embedding dives into a child of the middle term, compose
with the corresponding target child and dive; if both couple roots,
compose each matching pair of children. All subsidiary triples are
smaller. This proves the required quasi-order properties.

We need a finite-product fact. In any infinite sequence in a
well-quasi-order, some element lies below infinitely many later
elements. If this failed, each index would have only finitely many
later successors above it. Recursively choose a later index avoiding
the finitely many successors above all previously chosen indices.
The result would be an infinite bad subsequence. Applying the fact
repeatedly to the infinite set of successors above the chosen element
produces an infinite nondecreasing subsequence. For a finite tuple
of such orders, successively extract a nondecreasing subsequence in
each coordinate. Earlier coordinates retain their property. This
gives two tuples ordered coordinatewise, and proves that every finite
product is a well-quasi-order. The empty product has one element and
satisfies the same assertion.

Suppose now that the tree lemma fails. Choose a minimal bad sequence
\(t_0,t_1,\ldots\): after a chosen prefix, select a term of least
node size that permits an infinite bad continuation. Such a term
exists by the continuation hypothesis. Ties can be settled by a fixed
enumeration of terms over the finite alphabet. This is the usual
minimal-bad-sequence construction with a specified choice among
finite strings.

Let \(U\) be the set of all proper subterms of the \(t_i\).
We claim that \(U\) is well-quasi-ordered by embedding. Otherwise
take a bad sequence \(u_0,u_1,\ldots\) from it and associate to
each \(u_j\) its least parent index \(n_j\) such that it is a
proper subterm of \(t_{n_j}\). For a bounded set of parent indices
there are only finitely many such subterms. A bad sequence cannot
repeat a term, because embedding is reflexive. Consequently the
parent indices are unbounded on every tail, and we can extract a bad
subsequence with strictly increasing parent indices. Rename that
subsequence \(u_j\), with \(n_0<n_1<\cdots\).

The sequence
\(t_0,\ldots,t_{n_0-1},u_0,u_1,\ldots\) is bad. Its initial
segment was bad, and its tail is bad. If an initial term \(t_i\)
embedded into a tail term \(u_j\), then by transitivity and the
subterm embedding it would embed into \(t_{n_j}\), contradicting
\(i<n_0\leq n_j\) in the original bad sequence. But
\(\lvert u_0\rvert<\lvert t_{n_0}\rvert\), contradicting
minimality at position \(n_0\). Thus \(U\) is a well-quasi-order.

There are finitely many possible root symbols. Infinitely many of
the \(t_i\) therefore share one root \(f\) of fixed arity \(n\).
Their child tuples belong to \(U^n\). The finite-product fact gives
an earlier tuple embedding coordinatewise into a later one, and the
coupling clause then embeds the earlier whole term into the later
one. For \(n=0\) the terms are already identical. Both cases
contradict badness, proving the tree lemma.

For the ordering assertion, induct on the target tree in an embedding.
In a diving case, the source is equal to or smaller than a child by
induction, and the whole target is strictly greater than that child
by the subterm property. Transitivity gives strict inequality. In a
coupling case, every target child is equal to or greater than its
source child. Replacing unequal children one at a time by their smaller
partners and using context closure gives either literal equality of
the whole terms or a finite strict descending chain between them.
Transitivity gives the desired comparison.

If there were an infinite descending chain in the strict order, the
tree lemma would give indices \(i<j\) with \(t_i\unlhd t_j\).
The chain gives \(t_i>t_j\), whereas the embedding gives equality
or \(t_j>t_i\). Either contradicts transitivity and irreflexivity.
\end{proof}
\end{annexitem}

\begin{annexitem}{annex:p12.ro.lpo-termination}{thm:p12-ro-lpo-termination}{LPO well-foundedness, ground totality and rule termination}
\begin{proof}
In a descending LPO sequence, variable containment makes every
variable belong to the finite set occurring in its first term. Treat
those variables as distinct additional nullary symbols for embedding.
Together with the finite function signature this is a finite ranked
alphabet. The structural lemma and finite-tree lemma exclude an
infinite descending sequence. Context and substitution closure were
proved separately, so this is a reduction ordering.

When the precedence is total, prove comparability of distinct ground
terms \(s,t\) by induction on their combined size. If a child of
\(s\) is equal to or greater than \(t\), the subterm clause gives
\(s>_{\rm lpo}t\). The reverse case gives \(t>_{\rm lpo}s\).
If neither occurs, induction on child/whole pairs gives
\(t>_{\rm lpo}s_i\) for every child of \(s\), and
\(s>_{\rm lpo}t_j\) for every child of \(t\). Distinct roots
are ordered by the total precedence. For equal roots, the distinct
tuples have a first unequal coordinate at their common status, whose
entries are comparable by induction. Thus one of the root clauses
applies. Distinct nullary symbols are covered by the root precedence;
identical nullary symbols were excluded. This proves ground totality.

If every rule decreases, substitution and context closure make every
rewrite step decrease, and well-foundedness proves termination.
For the stated finite-TRS extension, every operator in any derivation
belongs to the union of the operators in the starting term and in
the finite rule set. Every variable in a reduct already occurred in
the original term, by the ordinary rule-variable condition. This is
again one fixed finite alphabet for the whole derivation. The
simplification-order argument excludes an infinite derivation even
without a global well-foundedness claim on the whole infinite signature.
\end{proof}
\end{annexitem}

\begin{annexitem}{annex:p12.ro.weights}{thm:p12-ro-weights}{Positive term weights and exact substitution calculations}
\begin{proof}
A variable has weight \(w_0\), and a nullary symbol has weight at
least \(w_0\). At a nonnullary symbol, induction gives at least
\(w_0\) for each child; adding the nonnegative symbol weight
therefore gives at least \(w_0\) for the whole term. This includes
zero-weight unary symbols and zero-weight branching symbols.

For a variable \(t=x\), the weight identity reads
\(w(\sigma(x))=w_0+(w(\sigma(x))-w_0)\); the occurrence identity
is equally immediate. For a nullary symbol both sums are empty and
substitution changes nothing. At a constructor, add the induction
identities for all children. Each occurrence of a variable belongs
to exactly one child, while the weight of the root symbol is unchanged.
The displayed sums therefore combine into precisely the identities
for the whole term. All sums are finite because the term is finite.
The first paragraph makes every increment
\(w(\sigma(x))-w_0\) nonnegative.
\end{proof}
\end{annexitem}

\begin{annexitem}{annex:p12.ro.kbo-structural}{thm:p12-ro-kbo-structural}{All KBO order, subterm and context cases}
\begin{proof}
The comparison is well-defined recursively: only the lexicographic
tie case calls it on proper arguments. It cannot compare a term
strictly with itself: the weight cannot be larger than itself, roots
cannot precede themselves, an identical tuple has no first unequal
coordinate, and \(z^k(x)\ne x\) for \(k>0\).
Every comparison implies weight nonincrease and variable-occurrence
dominance; at equal weights it implies equality of all variable counts.

A variable cannot be greater than a term, since its weight is minimal
and none of the equal-weight left-hand cases has a variable as source.
If a term has weight \(w_0\) and contains \(x\), it is
\(z^m(x)\) for some \(m\geq0\). Indeed, every leaf has weight
at least \(w_0\), so it has just that one variable leaf; every
operator above it must be unary of weight zero. Admissibility permits
only \(z\). Direct induction after canceling common initial \(z\)'s
shows that \(z^m(x)>_{\rm kb70}z^n(x)\) holds exactly when
\(m>n\). If no \(z\) exists only \(m=n=0\) is possible.

Prove transitivity on a triple \(s>_{\rm kb70}t>_{\rm kb70}u\)
by induction on its combined size. If either weight comparison is
strict, then \(w(s)>w(u)\), and transitivity of occurrence dominance
gives the first defining case. Otherwise all three weights and all
their variable-count vectors are equal. If \(u=x\), then
\(t=z^n(x)\), \(n>0\). The preceding observation makes
\(s=z^m(x)\), and \(s>t\) implies \(m>n\), hence \(s>u\).
If \(u\) is not a variable, all three terms have function-symbol
roots. If either root comparison is strict, transitivity of precedence
and equality in the other comparison give strict precedence between
the outer roots. If all three roots coincide, compare the first
unequal argument indices in the two lexicographic comparisons. The
earlier index decides when they differ. At a common index, induction
on the three proper arguments gives the required strict comparison.
All earlier arguments are equal, so the final lexicographic case
gives \(s>u\). Thus the order is transitive and irreflexive.

For an immediate subterm \(u\) of \(f(\bar t)\), occurrence
dominance always holds. If \(f\) has at least two arguments, the
other arguments contribute at least \(w_0>0\), giving a strict
weight decrease to \(u\). A positive-weight unary symbol gives
the same conclusion. For \(z(u)\), weights and variable counts are
equal. If \(u=x\), use the explicit unary-chain case. If the root
of \(u\) is not \(z\), use greatestness of \(z\). If
\(u=z(v)\), the equal-root lexicographic case reduces the desired
comparison to \(z(v)>v\), proved by induction on the smaller term
\(v\). There is no child case for constants. Transitivity extends
this to every proper subterm.

For one-argument context replacement, adding the unchanged root and
siblings adds equal quantities to the weights and to each variable
count. A strict weight comparison remains strict with dominance.
An equal-weight comparison retains equality of the whole variable-count
vectors, and the first unequal argument of the common root is the
changed one. The lexicographic clause applies there. Induction on
the context depth proves arbitrary context closure.

For distinct ground terms, variable conditions are vacuous. Unequal
weights decide. At equal weights, unequal roots are decided by total
precedence; equal roots have a first unequal argument, decided by
induction on combined term size. Identical constants are not distinct
terms. This proves totality on ground terms. Finally the occurrence
conditions immediately give the asserted containment of variable sets.
\end{proof}
\end{annexitem}

\begin{annexitem}{annex:p12.ro.kbo-substitution}{thm:p12-ro-kbo-substitution}{The exact historical variable conditions under substitution}
\begin{proof}
Induct on the combined size of a compared pair \(s>_{\rm kb70}t\).
Write \(d_x=\operatorname{occ}_x(s)-\operatorname{occ}_x(t)\).
In the strict-weight case every \(d_x\geq0\), and the weight
identity gives
\[
 w(s\sigma)-w(t\sigma)=w(s)-w(t)
       +\sum_x d_x\bigl(w(\sigma(x))-w_0\bigr)>0.
\]
The occurrence identity gives
\(\operatorname{occ}_y(s\sigma)-\operatorname{occ}_y(t\sigma)
=\sum_x d_x\operatorname{occ}_y(\sigma(x))\geq0\)
for every \(y\). Thus the first case persists.

In the equal-weight case every \(d_x=0\), so the two substituted
weights are equal and every substituted variable count is equal.
If the original comparison is \(z^k(x)>x\), the substituted source
is \(z^k(\sigma(x))\), which is greater than \(\sigma(x)\)
by the proper-subterm property. If it used unequal roots, those roots
are unaffected by substitution and preserve their precedence.
If it used equal roots with a first unequal argument, earlier argument
pairs remain equal and induction preserves the strict comparison of
that argument pair. Irreflexivity ensures that this pair has not become
equal. Hence the lexicographic clause persists. These exhaust the
original definition, including its zero-weight unary case.
\end{proof}
\end{annexitem}

\begin{annexitem}{annex:p12.ro.kbo-termination}{thm:p12-ro-kbo-termination}{Well-foundedness with zero-weight unary chains}
\begin{proof}
First consider ground terms. In a putative infinite descending chain
their natural-number weights cannot increase, so after finitely many
strict weight decreases there is an infinite tail of one weight \(W\).
For a finite ground tree let \(n_j\) count nodes of arity \(j\).
Counting its edges as both the number of nonroot nodes and the sum
of all arities gives
\[
 n_0=1+\sum_{j\geq2}(j-1)n_j.
\]
Since every nullary symbol has positive integer weight, \(n_0\leq W\).
Thus the total number of branching nodes is bounded by \(W-1\).
Every positive-weight unary node contributes at least one, so there
are at most \(W\) of them. If all unary weights are positive, term
size is therefore bounded at weight \(W\). The finite signature
then permits only finitely many such terms, contradicting a strict
descending chain.

If the exceptional unary symbol \(z\) has weight zero, erase all
its occurrences. The resulting core is still a ground term of weight
\(W\), and its size is bounded by the preceding count. There are
only finitely many possible cores. Take an infinite subsequence of
the descending chain with one fixed core; transitivity preserves
strict descent on this subsequence.

List the nodes of this core in left-to-right prefix order. Each
original term with that core is uniquely recovered by inserting a
nonnegative number of \(z\)'s immediately before each listed node.
It is therefore described by a fixed-length tuple of natural numbers.
For two terms with this core, KBO compares these tuples lexicographically.
To verify this, cancel common initial \(z\)'s. At the first unequal
insertion count, the side with more has root \(z\), whereas the
other has a core symbol different from \(z\). Their subtree weights
are equal, and greatestness of \(z\) makes the former larger.
If the insertion counts agree at a core root, its equal-root comparison
passes to the first differing child. Recursion follows precisely the
prefix order of the listed nodes. Thus no other comparison case can
change the first differing insertion count's decision.

Lexicographic order on a fixed finite tuple of natural numbers has
no infinite descending chain: its first coordinate can decrease only
finitely often and then stabilizes, after which induction on tuple
length applies to the tail coordinates. The empty tuple is trivial.
This contradicts the chosen fixed-core subsequence and proves ground
well-foundedness, including arbitrary lengths of \(z\)-chains.

For arbitrary terms, every variable in a descending sequence must
already occur in its first term, by variable containment. Choose a
nullary symbol of least positive weight and substitute it for those
finitely many variables. This is a finite-support substitution and
makes every term of the sequence ground. Substitution stability would turn
an infinite descending sequence into an infinite strictly descending
ground sequence. It cannot collapse a strict comparison to equality,
because strict comparison is preserved. Hence no descending sequence
of arbitrary terms exists. Combined with the proved context and
substitution properties, KBO is a reduction ordering. Decreasing
rewrite rules terminate by the usual rule-instance and context argument.
\end{proof}
\end{annexitem}

\begin{annexitem}{annex:p12.ro.examples}{thm:p12-ro-comparison-examples}{Complete orientation and failure calculations}
\begin{proof}
For the first Ackermann rule write
\(U=\mathsf A(\mathsf s(x),\mathsf s(y))\).
Proper-subterm comparisons give \(U>x\), \(U>y\) and
\(U>\mathsf s(x)\). Also
\(U>\mathsf A(\mathsf s(x),y)\): the first arguments agree,
the second decreases from \(\mathsf s(y)\) to \(y\), and the
source dominates both target arguments. Finally, comparison with
\(\mathsf A(x,\mathsf A(\mathsf s(x),y))\) succeeds because
its first argument decreases from \(\mathsf s(x)\) to \(x\),
and both target arguments have just been shown smaller than \(U\).

For \(U=\mathsf A(\mathsf s(x),\mathsf0)\), the target's first
argument \(x\) is smaller. Its other argument
\(\mathsf s(\mathsf0)\) is smaller by the strict root precedence
\(\mathsf A>_\Sigma\mathsf s\) and the subterm comparison
\(U>\mathsf0\). Thus the second rule decreases. For the third,
\(\mathsf A(\mathsf0,x)>\mathsf s(x)\) follows from that same
root precedence and domination of \(x\). The expanding false
variant has right side \(\mathsf A(x,U)\). Its first argument
alone would pass the lexicographic check, but domination would demand
\(U>U\), which is impossible. More generally the target has \(U\)
as a proper subterm, so a strict simplification ordering cannot orient
\(U\) toward it.

For the stated KBO data, every variable and \(e\) has weight one.
The comparison \(z(z(x))>x\) is the explicit chain case. For
\(m(e,x)>x\), weights are two and one and variable dominance holds.
Associativity has weight three on both sides with identical variable
counts; at their common root its first argument comparison is
\(m(x,y)>x\), decided by weights two and one with dominance.
For \(z(m(x,y))>m(z(y),z(x))\), both weights are two, both
variable-count vectors have one occurrence of each variable, and
the greatest root \(z\) exceeds \(m\). Thus every listed
orientation satisfies the precise 1970 definition.

If the greatestness requirement is removed, take a unary \(z\)
of weight zero, a unary \(g>_\Sigma z\) of weight one and a
constant \(c\) of weight one. The unchanged ground comparison
recipe gives \(g(c)>z(g(c))\) by equal weight and root precedence.
The equal-root argument comparison then gives
\(z^n(g(c))>z^{n+1}(g(c))\) for every \(n\geq0\).
This is an explicit infinite descent under the inadmissible data.

For a zero-weight constant failure, let \(c\) be nullary and
\(f\) binary, both of weight zero, with \(f>_\Sigma c\).
The raw ground weight/root/lexicographic recipe has
\(f(c,c)>c\). Put \(u=f(f(c,c),c)\) and \(H(t)=f(c,t)\).
Then \(u>H(u)\), since the first argument decreases from
\(f(c,c)\) to \(c\). Repeating the equal-root comparison gives
\(H^n(u)>H^{n+1}(u)\) for every \(n\), another infinite descent.
These examples are outside the admissible definition and explain its
restrictions; they are not counterexamples to the theorem.

If the variable condition were omitted, give unary \(f\) and
constant \(e\) weight one. The weight test alone would declare
\(f(x)>y\) for distinct variables, since two exceeds one.
Substitute \(e\) for \(x\) and \(f(f(e))\) for \(y\).
The resulting weights are two and three, reversing the proposed
ordering. The actual variable condition rejects the original comparison
because \(y\) is absent from its source.

For the LPO status comparison, put
\(S=m(m(x,y),u)\), \(T=m(x,m(y,u))\).
At left-to-right status, \(S>x\) and
\(S>m(y,u)\): the latter uses \(m(x,y)>y\) at the first
coordinate and domination of \(y,u\) by proper-subterm comparisons.
Then the root comparison of \(S,T\) uses \(m(x,y)>x\).
At right-to-left status, \(T>u\) and \(T>m(x,y)\): the latter
uses \(m(y,u)>y\) at the first coordinate inspected, with
domination of \(x,y\). The final comparison of \(T,S\) uses
\(m(y,u)>u\). This gives the two opposite orientations for two
different statuses, without contradicting strictness of either one.

Any substitution-stable strict order that oriented
\(m(x,y)>m(y,x)\) would, under the substitution interchanging
\(x,y\), also give the reverse strict comparison. Transitivity
and irreflexivity prohibit this. Both constructed orderings have those
properties, so neither can orient unrestricted commutativity.

For the final historical-variable example, both open terms have
weight two, but \(m(x,x)\) has two occurrences of \(x\) and
\(g(x)\) has one. Neither satisfies the equal-count tie condition,
and neither has strictly greater weight. They are therefore incomparable
under the historical definition. For a ground term \(t\), put
\(a=w(t)\geq1\). The two ground weights are \(2a\) and
\(a+1\). If \(a>1\), the first is larger. If \(a=1\), they
are equal and the root precedence \(m>_\Sigma g\) decides the
comparison. Ground terms have no variable-count obstruction. Thus
all ground instances decrease, showing that the open-term comparison
is a sufficient substitution-stable certificate, not a characterization
of every uniformly decreasing family of ground instances.
\end{proof}
\end{annexitem}
